\documentclass[../Odyssey_Project.tex]{subfiles}

\begin{document}
\setcounter{section}{7}
\section{Finite \texorpdfstring{$p$}{p}-groups}

We have seen that \nameref{thm:sylow} allows us to study finite groups by concentrating on their subgroups of prime-power order. This motivates the study of finite $p$-groups, a rich subject whose surface we shall only scratch in this section. Our first result is the most basic structural fact about finite $p$-groups.

\begin{theorem}
\label{thm:8.1}
If $P$ is a non-trivial finite $p$-group, then $Z(P)$ is also non-trivial. Moreover, if $1 \neq N \trianglelefteq P$, then $N \cap Z(P) \neq 1$.
\end{theorem}
\begin{proof}
Let $K_1=\{1\},K_2,\ldots,K_r$ be the conjugacy classes of $P$. Suppose first that $Z(P)=1$. If some $K_i$ with $i>1$ were a singleton, say $K_i=\{x\}$, then $gxg^{-1}=x$ for every $g\in P$, so $x\in Z(P)=1$, contradicting $i>1$. Hence $\lvert K_i\rvert>1$ for every $i>1$. Each $K_i$ is an orbit for the conjugation action of $P$ on itself, as in Proposition~\ref{prop:3.13}, and so $\lvert K_i\rvert$ divides $\lvert P\rvert$ by Corollary~\ref{cor:3.5}. Since $P$ is a $p$-group, it follows that $p$ divides $\lvert K_i\rvert$ for every $i>1$. Therefore
\[
    \lvert P\rvert
    =1+\lvert K_2\rvert+\cdots+\lvert K_r\rvert
    \equiv 1 \pmod p,
\]
which is impossible because $P$ is a non-trivial $p$-group. Thus $Z(P)\neq 1$.\\
Now let $1\neq N\trianglelefteq P$. Since $N$ is normal in $P$, it is a union of conjugacy classes of $P$, and since $1\in N$, this union contains $K_1=\{1\}$. If $N\cap Z(P)=1$, then every conjugacy class of $P$ contained in $N$ other than $K_1$ has size greater than $1$, and hence has size divisible by $p$ by the same orbit-size argument. Therefore $\lvert N\rvert\equiv 1\pmod p$, contradicting the fact that $N$ is a non-trivial subgroup of the $p$-group $P$. Hence $N\cap Z(P)\neq 1$.
\end{proof}

In particular, the only finite simple $p$-groups are the groups of prime order. Indeed, if $P$ is a non-trivial simple $p$-group, then Theorem~\ref{thm:8.1} gives $Z(P) \neq 1$. Since $Z(P) \trianglelefteq P$, simplicity forces $Z(P)=P$, so $P$ is abelian; hence $P$ has prime order.

\begin{lemma}
\label{lem:8.2}
If $G$ is a non-abelian group, then $G/Z(G)$ cannot be cyclic. In particular, if $P$ is a finite $p$-group, then $|P:Z(P)| \neq p$.
\end{lemma}
\begin{proof}
Suppose, for a contradiction, that $G/Z(G)$ is cyclic, say $G/Z(G)=\langle xZ(G)\rangle$ for some $x\in G$. Then every element of $G$ has the form $x^m z$ for some $m\in \Z$ and some $z\in Z(G)$. If $g=x^m z$ and $h=x^n w$, where $z,w\in Z(G)$, then
\[
    gh=x^m z x^n w=x^{m+n}zw=x^{n+m}wz=x^n w x^m z=hg.
\]
Thus any two elements of $G$ commute, so $G$ is abelian, contradicting the hypothesis. Therefore $G/Z(G)$ cannot be cyclic.\\
Now let $P$ be a finite $p$-group. If $P$ is abelian, then $Z(P)=P$, so $\lvert P:Z(P)\rvert=1$. If $P$ is non-abelian and $\lvert P:Z(P)\rvert=p$, then $P/Z(P)$ has prime order and is therefore cyclic, contradicting the first part. Hence $\lvert P:Z(P)\rvert\neq p$.
\end{proof}

Recall that a proper subgroup $H$ of a group $G$ is said to be \emph{maximal} if there is no proper subgroup of $G$ properly containing $H$. If $G$ is non-trivial and finite, then $G$ has maximal subgroups, and every proper subgroup of $G$ is contained in one. We now show that maximal subgroups of finite $p$-groups are especially well-behaved.

\begin{proposition}
\label{prop:8.3}
If $P$ is a finite $p$-group, then every maximal subgroup of $P$ is normal in $P$.
\end{proposition}
\begin{proof}
We argue by induction on $\lvert P\rvert$. If $\lvert P\rvert=1$, then $P$ has no maximal subgroups, and if $\lvert P\rvert=p$, then the only maximal subgroup is $1$, which is normal in $P$. Thus we may assume that $\lvert P\rvert>p$ and that the result holds for all finite $p$-groups of smaller order. Let $M$ be a maximal subgroup of $P$, and put $Z=Z(P)$. Since $Z$ is central, $MZ$ is a subgroup of $P$ containing $M$. By maximality, either $MZ=P$ or $MZ=M$.\\
If $MZ=P$, then every element of $P$ has the form $mz$ with $m\in M$ and $z\in Z$. For any $a\in M$, we have
\[
    (mz)a(mz)^{-1}=mza z^{-1}m^{-1}=mam^{-1}\in M,
\]
so $M\trianglelefteq P$.\\
It remains to consider the case $MZ=M$. Then $Z\leq M$. By Theorem~\ref{thm:8.1}, $Z\neq 1$, so $\lvert P/Z\rvert<\lvert P\rvert$. Moreover, by the \nameref{thm:correspondence}, $M/Z$ is a maximal subgroup of the finite $p$-group $P/Z$. The induction hypothesis gives $M/Z\trianglelefteq P/Z$, and applying the \nameref{thm:correspondence} once more gives $M\trianglelefteq P$.
\end{proof}

\begin{corollary}
\label{cor:8.4}
Any maximal subgroup of a finite $p$-group is of index $p$.
\end{corollary}
\begin{proof}
Let $M$ be a maximal subgroup of the finite $p$-group $P$. By Proposition~\ref{prop:8.3}, we have $M\trianglelefteq P$, so $P/M$ is a non-trivial finite $p$-group. By \nameref{thm:cauchy}, $P/M$ has a subgroup of order $p$. By the \nameref{thm:correspondence}, this subgroup has the form $L/M$ for some subgroup $L$ satisfying $M<L\leq P$. Since $M$ is maximal in $P$, we must have $L=P$. Therefore
\[
    \lvert P:M\rvert=\lvert P/M\rvert=\lvert L/M\rvert=p.
\]
\end{proof}

Our next task is to classify the finite abelian $p$-groups. The following lemma isolates the elementary fact about cyclic $p$-groups that will be used in the induction.

\begin{lemma}
\label{lem:8.5}
Every non-generator of a cyclic $p$-group $P$ is a $p$th power in $P$.
\end{lemma}
\begin{proof}
If $P=1$, then there is nothing to prove. Otherwise, write $P=\langle x\rangle$ with $\lvert P\rvert=p^a$ for some $a\geq 1$. By Theorem~\ref{thm:1.4}, the subgroup
\[
    Q=\langle x^p\rangle
\]
is the unique subgroup of $P$ of index $p$. Every element of $Q$ is a $p$th power in $P$, since
\[
    (x^p)^n=(x^n)^p
\]
for every $n\in \Z$.\\
We claim that $Q$ is precisely the set of non-generators of $P$. Since $Q<P$, no element of $Q$ can generate $P$. Conversely, if $y\in P$ does not generate $P$, then $\langle y\rangle$ is a proper subgroup of $P$, and hence is contained in some maximal subgroup $M$ of $P$. By Corollary~\ref{cor:8.4}, $M$ has index $p$, so by the uniqueness of $Q$ we have $M=Q$. Thus $y\in Q$, and therefore $y$ is a $p$th power in $P$.
\end{proof}

\begin{theorem}
\label{thm:8.6}
A finite abelian $p$-group is a direct product of cyclic $p$-groups.
\end{theorem}
\begin{proof}
We argue by induction on $\lvert P\rvert$. If $P=1$, there is nothing to prove, and if $\lvert P\rvert=p$, then $P$ is cyclic. Thus assume that $\lvert P\rvert>p$ and that the result holds for abelian $p$-groups of smaller order. Let $Q$ be a maximal subgroup of $P$. By Corollary~\ref{cor:8.4}, we have $\lvert P:Q\rvert=p$, and by induction we may write
\[
    Q=Q_1\times\cdots\times Q_s,
\]
where each $Q_i$ is cyclic of order $p^{a_i}$. Reordering if necessary, assume that $a_1\geq\cdots\geq a_s\geq 1$.\\
Choose $x\in P-Q$. Since $P/Q$ has order $p$, we have $x^p\in Q$. Write
\[
    x^p=y_1\cdots y_s,
\]
where $y_i\in Q_i$ for each $i$. If some $y_i$ is a non-generator of $Q_i$, then Lemma~\ref{lem:8.5} gives $y_i=z_i^p$ for some $z_i\in Q_i$. Replacing $x$ by $xz_i^{-1}$ keeps $x$ outside $Q$, and changes the $i$th coordinate of $x^p$ to $1$. Repeating this process, we may assume that each $y_i$ is either $1$ or a generator of $Q_i$.\\
If $x^p=1$, then $\langle x\rangle\cap Q=1$, because $xQ$ has order $p$ in $P/Q$. Hence $\langle x\rangle Q$ is a direct product of order $p\lvert Q\rvert=\lvert P\rvert$, so $P=\langle x\rangle\times Q$, and the result follows. We may therefore assume that $x^p\neq 1$. Let $j$ be minimal such that $y_j\neq 1$. Then $y_j$ is a generator of $Q_j$, so it has order $p^{a_j}$; since $a_j\geq a_i$ for all $i\geq j$, the element $x^p=y_j\cdots y_s$ has order $p^{a_j}$. Therefore $\langle x\rangle$ has order $p^{a_j+1}$.\\
Let $R$ be the direct product of all the $Q_i$ except $Q_j$. We claim that $\langle x\rangle\cap R=1$. Suppose $x^t\in R$, with $0\leq t<\lvert\langle x\rangle\rvert$. Since $xQ$ has order $p$ in $P/Q$, we must have $p\mid t$, say $t=mp$ with $0\leq m<p^{a_j}$. Then
\[
    x^t=(x^p)^m=y_j^m\cdots y_s^m.
\]
The $Q_j$-coordinate of this element is $y_j^m$. Since $y_j$ has order $p^{a_j}$ and $0\leq m<p^{a_j}$, this coordinate is the identity only when $m=0$. Thus $x^t\in R$ forces $t=0$, and so $\langle x\rangle\cap R=1$. It follows that $\langle x\rangle R$ is a direct product of order
\[
    \lvert\langle x\rangle\rvert\lvert R\rvert
    =p^{a_j+1}\frac{\lvert Q\rvert}{p^{a_j}}
    =p\lvert Q\rvert
    =\lvert P\rvert.
\]
Therefore $P=\langle x\rangle\times R$, which is a direct product of cyclic $p$-groups.
\end{proof}

We now consider how finite $p$-groups can be used to study arbitrary finite groups. The next result is a standard tool in finite group theory: it turns conjugacy of Sylow subgroups inside a normal subgroup into a factorization of the whole group.

\begin{namedtheorem}[Frattini Argument]
\label{thm:frattini-arg}
Let $G$ be a finite group, let $N \trianglelefteq G$, and let $P$ be a Sylow subgroup of $N$. Then $G = N_G(P)N$.
\end{namedtheorem}
\begin{proof}
Let $g\in G$. Since $N\trianglelefteq G$, we have $gNg^{-1}=N$, and hence $gPg^{-1}\leq N$. Moreover, $gPg^{-1}$ has the same order as $P$, so $gPg^{-1}$ is also a Sylow subgroup of $N$. By \nameref{thm:sylow}, applied inside $N$, there exists $n\in N$ such that
\[
    n(gPg^{-1})n^{-1}=P.
\]
Equivalently, $(ng)P(ng)^{-1}=P$, so $ng\in N_G(P)$. Thus $g=n^{-1}(ng)\in NN_G(P)$. Since $N\trianglelefteq G$, we have $NN_G(P)=N_G(P)N$, and hence $g\in N_G(P)N$. As $g$ was arbitrary, $G\leq N_G(P)N$. The reverse inclusion is immediate, since both $N_G(P)$ and $N$ are subgroups of $G$. Therefore $G=N_G(P)N$.
\end{proof}

\begin{theorem}
\label{thm:8.7}
The following statements about a finite group $G$ are equivalent:
\begin{enumerate}
    \item[(1)] Every Sylow subgroup of $G$ is normal in $G$.
    \item[(2)] $G$ is the direct product of its Sylow subgroups.
    \item[(3)] Every maximal subgroup of $G$ is normal in $G$.
\end{enumerate}
\end{theorem}
A finite group satisfying these equivalent conditions is said to be \emph{nilpotent}. Thus finite abelian groups and finite $p$-groups are nilpotent. A different but equivalent definition, one that extends beyond finite groups, will be discussed in Section~11.
\begin{proof}
We prove the implications in a cycle. First assume that every Sylow subgroup of $G$ is normal in $G$. Let $p_1,\ldots,p_n$ be the distinct prime divisors of $\lvert G\rvert$, and let $P_i$ be the Sylow $p_i$-subgroup of $G$. By hypothesis, each $P_i$ is normal in $G$. We claim by induction that
\[
    P_1\cdots P_i
\]
is a normal subgroup of $G$ of order $\lvert P_1\rvert\cdots\lvert P_i\rvert$ for every $i$. This is clear for $i=1$. If it holds for $i$, then Proposition~\ref{prop:1.7} gives $P_1\cdots P_iP_{i+1}\trianglelefteq G$, and
\[
    (P_1\cdots P_i)\cap P_{i+1}=1
\]
because its order divides both $\lvert P_1\rvert\cdots\lvert P_i\rvert$ and $\lvert P_{i+1}\rvert$, which are relatively prime. Hence Proposition~\ref{prop:3.12} gives
\[
    \lvert P_1\cdots P_iP_{i+1}\rvert
    =\lvert P_1\cdots P_i\rvert\lvert P_{i+1}\rvert
    =\lvert P_1\rvert\cdots\lvert P_{i+1}\rvert.
\]
The claim follows. Taking $i=n$, we get $\lvert P_1\cdots P_n\rvert=\lvert G\rvert$, so $G=P_1\cdots P_n$. Moreover, for each $i$, the subgroup $P_i\cap P_1\cdots P_{i-1}P_{i+1}\cdots P_n$ has order dividing $\lvert P_i\rvert$ and also dividing the product of the other Sylow orders, so it is trivial. Therefore $G$ is the internal direct product of its Sylow subgroups. This proves (1) $\Rightarrow$ (2).\\
Now assume that $G=P_1\times\cdots\times P_n$, where the $P_i$ are the Sylow subgroups of $G$. Let $M$ be a maximal subgroup of $G$. Since the orders $\lvert P_i\rvert$ are pairwise relatively prime, Proposition~\ref{prop:2.10} gives
\[
    M=(M\cap P_1)\times\cdots\times(M\cap P_n).
\]
At least one factor $M\cap P_j$ is proper in $P_j$, since otherwise $M=G$. Let $j$ be any index for which $M\cap P_j<P_j$. If $M\cap P_j$ were not maximal in $P_j$, then replacing it by an intermediate subgroup of $P_j$ would give a subgroup strictly between $M$ and $G$, contradicting the maximality of $M$. Hence each proper factor $M\cap P_j$ is maximal in $P_j$. There cannot be two such proper factors: if $M\cap P_j<P_j$ and $M\cap P_k<P_k$ with $j\neq k$, then replacing $M\cap P_j$ by $P_j$ while leaving the $k$th factor proper again gives a subgroup strictly between $M$ and $G$. Thus there is exactly one index $j$ for which $M\cap P_j<P_j$, and this subgroup is maximal in the finite $p_j$-group $P_j$. By Proposition~\ref{prop:8.3}, $M\cap P_j\trianglelefteq P_j$. The other factors are the whole $P_i$, so they are normal in $P_i$. It follows directly from the direct product multiplication that $M\trianglelefteq G$. This proves (2) $\Rightarrow$ (3).\\
Finally assume that every maximal subgroup of $G$ is normal in $G$. Let $P$ be a Sylow subgroup of $G$. Suppose, for a contradiction, that $P$ is not normal in $G$. Then $N_G(P)<G$, so choose a maximal subgroup $M$ of $G$ with
\[
    P\leq N_G(P)\leq M<G.
\]
By hypothesis, $M\trianglelefteq G$. Also $P$ is a Sylow subgroup of $M$, since $P\leq M$ and $P$ already has the full $p$-part possible in $G$. Applying the \nameref{thm:frattini-arg} to the normal subgroup $M\trianglelefteq G$ and its Sylow subgroup $P$, we obtain
\[
    G=N_G(P)M.
\]
But $N_G(P)\leq M$, so the right hand side is just $M$, contradicting $M<G$. Therefore every Sylow subgroup of $G$ is normal in $G$, proving (3) $\Rightarrow$ (1).
\end{proof}

As a consequence of this result and our previous classification of finite abelian $p$-groups, we obtain the usual basis theorem for finite abelian groups.

\begin{corollary}
\label{cor:8.8}
A finite abelian group is a direct product of cyclic $p$-groups.
\end{corollary}
\begin{proof}
Let $G$ be a finite abelian group. Since every subgroup of an abelian group is normal, every Sylow subgroup of $G$ is normal in $G$. By Theorem~\ref{thm:8.7}, $G$ is the direct product of its Sylow subgroups. Each Sylow subgroup is a finite abelian $p$-group, and hence is a direct product of cyclic $p$-groups by Theorem~\ref{thm:8.6}. Combining these direct-product decompositions gives the result.
\end{proof}

We close this section with a discussion of $p$-groups of small order. Any group of order $p$ is isomorphic with $\Z_p$. By Theorem~\ref{thm:8.1} and Lemma~\ref{lem:8.2}, any group of order $p^2$ is abelian, and hence is isomorphic with either $\Z_{p^2}$ or $\Z_p \times \Z_p$ by Theorem~\ref{thm:8.6}. The same theorem shows that the abelian groups of order $p^3$ are precisely $\Z_{p^3}$, $\Z_{p^2} \times \Z_p$, and $\Z_p \times \Z_p \times \Z_p$. We therefore turn to the non-abelian groups of order $p^3$. For odd primes $p$, the following commutator identities provide the technical input for the classification.

\begin{lemma}
\label{lem:8.9}
Let $G$ be a group and let $x, y \in G$. If $[x,y] \in Z(G)$, then $[x^n, y] = [x,y]^n$ and $x^n y^n = (xy)^n [x,y]^{\frac{n(n-1)}{2}}$ for any $n \in \N$.
\end{lemma}
\begin{proof}
Put $c=[x,y]$. Since $c\in Z(G)$, the relation $xy=cyx$ allows us to move $y$ past powers of $x$. We first prove by induction on $n$ that
\[
    x^n y=c^n yx^n.
\]
The case $n=1$ is exactly $xy=cyx$. If $x^n y=c^n yx^n$, then
\[
    x^{n+1}y=x(c^n yx^n)=c^n xyx^n=c^{n+1}yx^{n+1},
\]
because $c$ is central and $xy=cyx$. Therefore $x^n y=c^n yx^n$ for all $n\in \N$, and hence
\[
    [x^n,y]=x^n yx^{-n}y^{-1}=c^n=[x,y]^n.
\]
It remains to prove the second identity. We again argue by induction on $n$. The case $n=1$ is immediate. Suppose that
\[
    x^n y^n=(xy)^n c^{\frac{n(n-1)}{2}}.
\]
Using the identity already proved, we have $x^n y=[x^n,y]yx^n=c^n yx^n$. Thus
\[
    x^{n+1}y^{n+1}
    =x(x^n y)y^n
    =xc^n yx^n y^n
    =c^n xyx^n y^n.
\]
By the induction hypothesis, this becomes
\[
    x^{n+1}y^{n+1}
    =c^n(xy)(xy)^n c^{\frac{n(n-1)}{2}}
    =(xy)^{n+1}c^{\frac{n(n-1)}{2}+n}
    =(xy)^{n+1}c^{\frac{(n+1)n}{2}}.
\]
Substituting back $c=[x,y]$ gives the desired formula.
\end{proof}

\begin{proposition}
\label{prop:8.10}
If $p$ is an odd prime, then there is exactly one isomorphism class of non-abelian groups of order $p^3$ having elements of order $p^2$.
\end{proposition}
\begin{proof}
Let $P$ be a non-abelian group of order $p^3$ having an element of order $p^2$. We first show that $P$ can be written as a semidirect product of $\Z_{p^2}$ by $\Z_p$. Choose $y\in P$ with $o(y)=p^2$, and choose $x\in P-\langle y\rangle$. Since $\langle y\rangle$ has index $p$, it is maximal in $P$, and hence $\langle y\rangle\trianglelefteq P$ by Proposition~\ref{prop:8.3}.\\
We may arrange that $x$ has order $p$. If $x^p=1$, there is nothing to change, so suppose that $x^p\neq 1$. Since $P$ is non-abelian of order $p^3$, Theorem~\ref{thm:8.1} and Lemma~\ref{lem:8.2} force $\lvert Z(P)\rvert=p$. Also $\langle y\rangle\cap Z(P)\neq 1$ by Theorem~\ref{thm:8.1}, so the unique subgroup of order $p$ in the cyclic group $\langle y\rangle$ gives
\[
    Z(P)=\langle y^p\rangle.
\]
The quotient $P/Z(P)$ has order $p^2$, and it is not cyclic by Lemma~\ref{lem:8.2}; hence every element of $P/Z(P)$ has order dividing $p$. Thus $x^p\in Z(P)$. Since $x^p\neq 1$, we have $x^p=y^{kp}$ for some $1\leq k<p$. Replacing $y$ by $y^{-k}$, which is still a generator of $\langle y\rangle$, we may assume that
\[
    x^p=y^{-p}.
\]
Now $P/Z(P)$ is abelian, so Proposition~\ref{prop:2.6} gives $P'\leq Z(P)$. In particular, $[x,y]\in Z(P)$. By Lemma~\ref{lem:8.9},
\[
    x^p y^p=(xy)^p[x,y]^{\frac{p(p-1)}{2}}.
\]
The left hand side is $1$, and $[x,y]^{\frac{p(p-1)}{2}}=1$, since $[x,y]\in Z(P)$ has order dividing $p$ and $p$ divides $\frac{p(p-1)}{2}$ for odd $p$. Therefore $(xy)^p=1$. Also $xy\notin\langle y\rangle$, because $x\notin\langle y\rangle$. Replacing $x$ by $xy$, we have $o(x)=p$, $o(y)=p^2$, and $x\notin\langle y\rangle$.\\
It follows that $\langle x\rangle\cap\langle y\rangle=1$, and hence
\[
    \lvert \langle y\rangle\langle x\rangle\rvert
    =p^2p
    =p^3.
\]
Thus $P=\langle y\rangle\langle x\rangle$, with $\langle y\rangle\trianglelefteq P$, so $P$ is a semidirect product of $\langle y\rangle\cong \Z_{p^2}$ by $\langle x\rangle\cong \Z_p$. Since $P$ is non-abelian, the corresponding homomorphism $\varphi\colon \Z_p\to \Aut(\Z_{p^2})$ is non-trivial, and hence is a monomorphism.\\
It remains to see that this semidirect product is unique up to isomorphism. By Proposition~\ref{prop:2.1}, $\Aut(\Z_{p^2})\cong(\Z/p^2\Z)^{\times}$, an abelian group of order $p(p-1)$. By \nameref{thm:sylow}, $\Aut(\Z_{p^2})$ has a subgroup of order $p$, and because $\Aut(\Z_{p^2})$ is abelian, this Sylow $p$-subgroup is unique. Therefore every monomorphism $\Z_p\to \Aut(\Z_{p^2})$ has the same image. Proposition~\ref{prop:2.11} then implies that all non-trivial semidirect products of $\Z_{p^2}$ by $\Z_p$ are isomorphic. Such a semidirect product exists by choosing an isomorphism from $\Z_p$ onto the unique subgroup of order $p$ in $\Aut(\Z_{p^2})$, and it is non-abelian because the action is non-trivial. Hence there is exactly one isomorphism class of non-abelian groups of order $p^3$ having an element of order $p^2$.
\end{proof}

\begin{proposition}
\label{prop:8.11}
If $p$ is an odd prime, then there is exactly one isomorphism class of non-abelian groups of order $p^3$ having no elements of order $p^2$.
\end{proposition}
\begin{proof}
Let $P$ be a non-abelian group of order $p^3$ having no elements of order $p^2$. Choose a maximal subgroup $M$ of $P$. By Corollary~\ref{cor:8.4}, we have $\lvert P:M\rvert=p$, and by Proposition~\ref{prop:8.3}, we have $M\trianglelefteq P$. Thus $\lvert M\rvert=p^2$. By Lemma~\ref{lem:8.2}, every group of order $p^2$ is abelian; since $P$ has no elements of order $p^2$, the subgroup $M$ is not cyclic, and hence Theorem~\ref{thm:8.6} gives $M\cong \Z_p\times \Z_p$. Choose $x\in P-M$. Again using the hypothesis, $x$ has order $p$. Also $\langle x\rangle\cap M=1$, and so
\[
    \lvert M\langle x\rangle\rvert
    =\lvert M\rvert\lvert\langle x\rangle\rvert
    =p^3.
\]
Therefore $P=M\langle x\rangle$, and $P$ is a semidirect product of $M\cong \Z_p\times\Z_p$ by $\langle x\rangle\cong \Z_p$. Since $P$ is non-abelian, the corresponding homomorphism $\varphi\colon \Z_p\to\Aut(\Z_p\times\Z_p)$ is non-trivial, and hence is a monomorphism.\\
We now show that all such semidirect products are isomorphic. Put $E=\Z_p\times\Z_p$. By Proposition~\ref{prop:4.1}, $\Aut(E)\cong \GL(2,p)$, and by Proposition~\ref{prop:4.2},
\[
    \lvert \Aut(E)\rvert
    =\lvert \GL(2,p)\rvert
    =(p^2-1)(p^2-p)
    =p(p-1)^2(p+1).
\]
Thus a Sylow $p$-subgroup of $\Aut(E)$ has order $p$, so all subgroups of order $p$ in $\Aut(E)$ are conjugate by \nameref{thm:sylow}. If $\varphi,\psi\colon \Z_p\to\Aut(E)$ are monomorphisms, then there is some $f\in\Aut(E)$ such that
\[
    \psi(\Z_p)=f\varphi(\Z_p)f^{-1}.
\]
Equivalently, if $\widehat f$ denotes conjugation by $f$ on $\Aut(E)$, then $\psi(\Z_p)=(\widehat f\circ\varphi)(\Z_p)$. By Proposition~\ref{prop:2.12},
\[
    E\rtimes_{\widehat f\circ\varphi}\Z_p\cong E\rtimes_{\varphi}\Z_p,
\]
while Proposition~\ref{prop:2.11} gives
\[
    E\rtimes_{\widehat f\circ\varphi}\Z_p\cong E\rtimes_{\psi}\Z_p,
\]
because $\Z_p$ is cyclic and the two actions have the same image. Hence all non-trivial semidirect products of $\Z_p\times\Z_p$ by $\Z_p$ are isomorphic.\\
It remains only to check that this isomorphism class actually occurs among groups with no elements of order $p^2$. Let $E$ be written additively as a two-dimensional vector space over $\mathbb{F}_p$, and let a generator $t$ of $\Z_p$ act on $E$ by
\[
    A=
    \begin{pmatrix}
        1 & 1\\
        0 & 1
    \end{pmatrix}.
\]
Writing $A=I+N$, we have $N^2=0$, so $A^p=I+pN=I$ and $A\neq I$. Thus this defines a non-trivial semidirect product $E\rtimes\langle t\rangle$. We claim that every element has order dividing $p$. An element has the form $(v,t^i)$. If $i=0$, then $(v,1)^p=(0,1)$ because $E$ has exponent $p$. If $i\neq 0$, then
\[
    (v,t^i)^p
    =
    \left(\sum_{j=0}^{p-1}A^{ij}v,1\right)
    =
    \left(\sum_{j=0}^{p-1}(I+ijN)v,1\right)
    =
    \left(\left(pI+i\frac{p(p-1)}{2}N\right)v,1\right)
    =(0,1),
\]
since we are working over $\mathbb{F}_p$ and $p$ is odd. Therefore this semidirect product is non-abelian of order $p^3$ and has no elements of order $p^2$. Combined with the uniqueness just proved, this gives exactly one isomorphism class.
\end{proof}

It remains to mention the case $p=2$. A non-abelian group of order $8$ has an element of order $4$ by Exercise~\ref{ex:1.2}. If it also has an element of order $2$ outside the cyclic subgroup of order $4$, then it is the non-trivial semidirect product $\Z_4 \rtimes \Z_2$, namely the dihedral group $D_8$. This will happen unless the group has a unique element $t$ of order $2$. In that remaining case, choose an element $y$ of order $4$, put $Q=\langle y\rangle$, and choose $x$ so that $P/Q=\langle xQ\rangle$. Since $x$ cannot have order $2$, we must have $x^2=t=y^2$. Every element can then be written uniquely as $x^a y^b$, where $0\leq a\leq 3$ and $0\leq b\leq 1$, and the multiplication is forced by the relation $yx=x^3y$. Thus, if this case exists, it is unique up to isomorphism.\\
To see that it exists, work in $\SL(2,3)$ over the field $\mathbb{F}_3$ and take
\[
    x=
    \begin{pmatrix}
        1 & 1\\
        1 & 2
    \end{pmatrix},
    \qquad
    y=
    \begin{pmatrix}
        0 & 2\\
        1 & 0
    \end{pmatrix}.
\]
The subgroup generated by these two matrices is non-abelian of order $8$, and it has the unique element
\[
    \begin{pmatrix}
        2 & 0\\
        0 & 2
    \end{pmatrix}
\]
of order $2$. This group is the quaternion group $Q_8$, which appeared in Exercise~\ref{ex:2.11} and cannot be written non-trivially as a semidirect product.

\subsection*{Exercises}

\begin{exercise}
\label{ex:8.1}
Show that if $P$ is a non-cyclic finite $p$-group, then $P$ has a normal subgroup $N$ such that $P/N \cong \Z_p \times \Z_p$.
\end{exercise}
\begin{solution}
Suppose first that $P$ has only one maximal subgroup, say $M$. If $x\in P-M$, then $\langle x\rangle$ is not contained in $M$. Since every proper subgroup of the finite group $P$ is contained in a maximal subgroup, $\langle x\rangle$ cannot be proper. Hence $\langle x\rangle=P$, so $P$ is cyclic. Therefore, since $P$ is non-cyclic, it has at least two distinct maximal subgroups. Choose two of them, say $M$ and $K$.\\
By Proposition~\ref{prop:8.3}, both $M$ and $K$ are normal in $P$, and by Corollary~\ref{cor:8.4} both have index $p$. Put
\[
    N=M\cap K.
\]
Then $N\trianglelefteq P$. Since $M\neq K$, we have $K<MK\leq P$, and the maximality of $K$ gives $MK=P$. Hence
\[
    \lvert P\rvert=\lvert MK\rvert
    =\frac{\lvert M\rvert\lvert K\rvert}{\lvert M\cap K\rvert}
    =\frac{(\lvert P\rvert/p)(\lvert P\rvert/p)}{\lvert N\rvert},
\]
so $\lvert P:N\rvert=p^2$. Thus $P/N$ has order $p^2$. Moreover, $M/N$ and $K/N$ are distinct subgroups of $P/N$ of order $p$. A cyclic group of order $p^2$ has a unique subgroup of order $p$, so $P/N$ is not cyclic. By Lemma~\ref{lem:8.2}, every group of order $p^2$ is abelian, and then Theorem~\ref{thm:8.6} gives $P/N\cong \Z_p\times \Z_p$.
\end{solution}

\begin{exercise}
\label{ex:8.2}
Let $P$ be a group of order $p^n$. Show that $P$ has a normal subgroup $N_a$ of order $p^a$ for every $0 \leq a \leq n$, and that these subgroups can be chosen so that $N_a$ is contained in $N_b$ whenever $a \leq b$.
\end{exercise}
\begin{solution}
We argue by induction on $n$. If $n=0$, then $P=1$, and we take $N_0=1$. Now assume that $n\geq 1$ and that the result is known for groups of order $p^{n-1}$. By Theorem~\ref{thm:8.1}, the center $Z(P)$ is non-trivial. Hence, by \nameref{thm:cauchy}, there is a subgroup
\[
    C\leq Z(P)
\]
of order $p$. Since $C$ is central, it is normal in $P$.\\
Let $\pi\colon P\to P/C$ be the quotient map. The group $P/C$ has order $p^{n-1}$, so by the induction hypothesis there are normal subgroups
\[
    \overline{N}_0\leq \overline{N}_1\leq \cdots \leq \overline{N}_{n-1}
\]
of $P/C$, with $\lvert \overline{N}_a\rvert=p^a$ for each $0\leq a\leq n-1$. Define
\[
    N_0=1,
    \qquad
    N_a=\pi^{-1}(\overline{N}_{a-1})
    \quad\text{for }1\leq a\leq n.
\]
Each $N_a$ is normal in $P$, because it is the preimage of a normal subgroup under the homomorphism $\pi$. Also $C\leq N_a$ for $a\geq 1$, and
\[
    N_a/C\cong \overline{N}_{a-1}.
\]
Therefore
\[
    \lvert N_a\rvert
    =\lvert C\rvert\,\lvert \overline{N}_{a-1}\rvert
    =p\cdot p^{a-1}
    =p^a
\]
for every $1\leq a\leq n$, while $\lvert N_0\rvert=1=p^0$. Finally, the inclusions among the $\overline{N}_a$ give
\[
    N_0\leq N_1\leq \cdots \leq N_n,
\]
so the subgroups may be chosen with $N_a\leq N_b$ whenever $a\leq b$.
\end{solution}

\begin{exercise}
\label{ex:8.3}
Let $G = \mathrm{GL}(n,p)$, and let $P$ be a Sylow $p$-subgroup of $G$. What is the order of $Z(P)$? What is the order of $Z(P/Z(P))$? If we let $Z_2(P) \leq P$ be such that $Z_2(P)/Z(P) = Z(P/Z(P))$ and continue in this way, what happens?
\end{exercise}
\begin{solution}
Let $U$ be the subgroup of upper unitriangular matrices in $\mathrm{GL}(n,p)$. Then $\lvert U\rvert=p^{n(n-1)/2}$, which is the $p$-part of $\lvert \mathrm{GL}(n,p)\rvert$, so we may take $P=U$. If $n=1$, then $P=1$ and there is nothing to prove, so assume $n\geq 2$. For $i<j$, put
\[
    X_{ij}=\{I+aE_{ij}:a\in\mathbb{F}_p\},
\]
where $E_{ij}$ is the matrix with $1$ in the $(i,j)$-position and $0$ elsewhere. These are the root subgroups of $U$. The basic matrix-unit calculation shows that a commutator between $X_{ij}$ and $X_{kl}$ can be non-trivial only when $j=k$ or $l=i$.\\
The subgroup $X_{1n}$ commutes with every root subgroup: the conditions $n=k$ and $l=1$ are impossible for $k<l$ inside $U$. Conversely, in an arbitrary element of $U$, any non-zero $X_{ij}$-coordinate with $(i,j)\neq(1,n)$ is detected by an adjacent root subgroup: either $i>1$ and $X_{i-1,i}$ detects it, or $j<n$ and $X_{j,j+1}$ detects it. Thus centrality forces all root coordinates except possibly the $X_{1n}$-coordinate to vanish, and so
\[
    Z(P)=X_{1n}
\]
and $\lvert Z(P)\rvert=p$.\\
If $n=2$, then $P$ is already abelian, so the quotient case is trivial. Assume now that $n\geq 3$. Let $Z_2(P)\leq P$ be defined by $Z_2(P)/Z(P)=Z(P/Z(P))$. Modulo $Z(P)=X_{1n}$, a root subgroup becomes central exactly when all of its possible non-trivial commutators land in $X_{1n}$. This happens precisely for
\[
    X_{1,n-1},\qquad X_{1n},\qquad X_{2n}.
\]
Hence
\[
    Z_2(P)=X_{1,n-1}X_{1n}X_{2n},
\]
so $\lvert Z_2(P)\rvert=p^3$, and therefore
\[
    \lvert Z(P/Z(P))\rvert
    =\lvert Z_2(P)/Z(P)\rvert
    =\frac{p^3}{p}
    =p^2.
\]
Continuing in this way gives the upper central series of $P$. If $Z_1(P)=Z(P)$, then for $1\leq r\leq n-1$,
\[
    Z_r(P)=\langle X_{ij}:j-i\geq n-r\rangle.
\]
Equivalently, $Z_r(P)$ consists of the upper unitriangular matrices whose non-zero off-diagonal entries can occur only on the $r$ outer superdiagonals nearest the top-right corner. Thus each step adds the next outer superdiagonal, and
\[
    \lvert Z_r(P)\rvert=p^{1+2+\cdots+r}=p^{r(r+1)/2}
\]
for $1\leq r\leq n-1$. After $n-1$ steps, we reach $Z_{n-1}(P)=P$.
\end{solution}

\begin{exercise}
\label{ex:8.4}
Let $U$ be the subgroup of $\mathrm{GL}(n,p)$ consisting of the upper unitriangular matrices, and let $Q$ be the subgroup of $U$ consisting of matrices whose $(i,j)$-entry is zero whenever $1 < i < j < n$. Determine $Z(Q)$, and show that $Q/Z(Q)$ is abelian.
\end{exercise}
\begin{solution}
For $i<j$, write
\[
    X_{ij}=\{I+aE_{ij}:a\in\mathbb{F}_p\}.
\]
The subgroup $Q$ consists of the upper unitriangular matrices whose possible non-zero off-diagonal entries lie either in the first row or in the last column. Equivalently,
\[
    Q=\langle X_{1j},X_{in}:2\leq j\leq n,\ 1\leq i<n\rangle.
\]
The only possible non-trivial commutators among these root subgroups occur when a first-row root meets the corresponding last-column root:
\[
    [X_{1j},X_{jn}]\leq X_{1n}
    \qquad (2\leq j\leq n-1).
\]
All other commutators between the displayed root subgroups are trivial. In particular, $X_{1n}\leq Z(Q)$.\\
We now show that there are no other central root directions. Let $g\in Q$. If the $X_{1j}$-coordinate of $g$ is non-zero for some $2\leq j\leq n-1$, then commuting $g$ with an element of $X_{jn}$ gives a non-trivial element of $X_{1n}$, so $g$ is not central. Similarly, if the $X_{jn}$-coordinate of $g$ is non-zero for some $2\leq j\leq n-1$, then commuting $g$ with an element of $X_{1j}$ gives a non-trivial element of $X_{1n}$. Thus a central element of $Q$ can have only its $X_{1n}$-coordinate non-zero, and hence
\[
    Z(Q)=X_{1n}.
\]
Finally, since every commutator in $Q$ lies in $X_{1n}=Z(Q)$, all commutators become trivial in the quotient $Q/Z(Q)$. Therefore $Q/Z(Q)$ is abelian.
\end{solution}

\begin{exercise}
\label{ex:8.5}
Show that subgroups and quotient groups of finite nilpotent groups are nilpotent, and that direct products of finite nilpotent groups are nilpotent.
\end{exercise}
\begin{solution}
Let $G$ be a finite nilpotent group. By Theorem~\ref{thm:8.7}, we may write
\[
    G=P_1\times\cdots\times P_r,
\]
where $P_i$ is the Sylow $p_i$-subgroup of $G$.\\
First let $H\leq G$. We claim that
\[
    H=(H\cap P_1)\times\cdots\times(H\cap P_r).
\]
Indeed, if $h\in H$, write $h=x_1\cdots x_r$ with $x_i\in P_i$. Since the factors have pairwise relatively prime orders, each $x_i$ is a power of $h$, and hence $x_i\in H$. Thus $x_i\in H\cap P_i$, proving the claim. Each $H\cap P_i$ is a normal subgroup of $H$, because $P_i\trianglelefteq G$. These intersections are precisely the Sylow subgroups of $H$ which are non-trivial. Hence every Sylow subgroup of $H$ is normal in $H$, so $H$ is nilpotent by Theorem~\ref{thm:8.7}.\\
Now let $N\trianglelefteq G$. For each $i$, the subgroup
\[
    P_iN/N
\]
is a normal $p_i$-subgroup of $G/N$, and
\[
    G/N=(P_1N/N)\cdots(P_rN/N).
\]
These factors have pairwise relatively prime orders and commute with one another, since the $P_i$ commute in $G$. Therefore the Sylow subgroups of $G/N$ are normal, and $G/N$ is nilpotent.\\
Finally, let $G_1,\ldots,G_m$ be finite nilpotent groups. If $S_{i,p}$ denotes the Sylow $p$-subgroup of $G_i$, with $S_{i,p}=1$ when $p\nmid \lvert G_i\rvert$, then
\[
    S_{1,p}\times\cdots\times S_{m,p}
\]
is the Sylow $p$-subgroup of $G_1\times\cdots\times G_m$. Each $S_{i,p}$ is normal in $G_i$, so this product is normal in $G_1\times\cdots\times G_m$. Thus every Sylow subgroup of the direct product is normal, and the direct product is nilpotent.
\end{solution}

\begin{exercise}
\label{ex:8.6}
Let $U$ be the subgroup of $\mathrm{GL}(3,p)$ consisting of the upper unitriangular matrices. Show that if $p$ is an odd prime, then $U$ is a non-abelian group of order $p^3$ having no elements of order $p^2$. If $p = 2$, with which group of order $8$ is $U$ isomorphic?
\end{exercise}
\begin{solution}
Write an element of $U$ as
\[
    u(a,b,c)=
    \begin{pmatrix}
        1 & a & c\\
        0 & 1 & b\\
        0 & 0 & 1
    \end{pmatrix},
    \qquad a,b,c\in\mathbb{F}_p.
\]
There are $p$ choices for each of $a,b,c$, so $\lvert U\rvert=p^3$. A direct multiplication gives
\[
    u(a,b,c)u(a',b',c')
    =
    u(a+a',b+b',c+c'+ab').
\]
Thus $U$ is not abelian, since
\[
    u(1,0,0)u(0,1,0)=u(1,1,1),
    \qquad
    u(0,1,0)u(1,0,0)=u(1,1,0).
\]
Now let
\[
    N=aE_{12}+bE_{23}+cE_{13}.
\]
Then $u(a,b,c)=I+N$, while $N^2=abE_{13}$ and $N^3=0$. Hence, for $m\geq 1$,
\[
    u(a,b,c)^m
    =(I+N)^m
    =I+mN+\binom{m}{2}N^2
    =
    u\left(ma,mb,mc+\binom{m}{2}ab\right).
\]
If $p$ is odd, then in $\mathbb{F}_p$ we have $pa=pb=pc=0$ and
\[
    \binom{p}{2}=\frac{p(p-1)}{2}\equiv 0 \pmod p.
\]
Therefore $u(a,b,c)^p=I$ for every element of $U$. Thus every element has order dividing $p$, so no element can have order $p^2$.

It remains to consider $p=2$. Put
\[
    r=u(1,1,0),
    \qquad
    s=u(1,0,0).
\]
Then $r^2=u(0,0,1)$, so $r$ has order $4$, while $s$ has order $2$. Also
\[
    srs^{-1}=srs=u(1,1,1)=r^{-1}.
\]
Since $s\notin\langle r\rangle$, the elements $r$ and $s$ generate all of $U$. Thus $U$ has the usual presentation of the dihedral group of order $8$, and so
\[
    U\cong D_8.
\]
\end{solution}

\subsection*{Further Exercises}

\begin{exercise}
\label{ex:8.7}
Show that a finite group $G$ has a largest nilpotent normal subgroup, in the sense that it contains all nilpotent normal subgroups of $G$. (This subgroup is called the \emph{Fitting subgroup} of $G$.)
\end{exercise}
\begin{solution}
We first prove that the product of two nilpotent normal subgroups of $G$ is again nilpotent. Let $A,B\trianglelefteq G$ be nilpotent. By Theorem~\ref{thm:8.7}, write
\[
    A=\prod_p A_p,
    \qquad
    B=\prod_p B_p,
\]
where $A_p$ and $B_p$ are the Sylow $p$-subgroups of $A$ and $B$, with the trivial subgroup inserted when $p$ does not divide the relevant order. Since $A$ and $B$ are nilpotent, these Sylow subgroups are unique in $A$ and $B$, respectively. Hence each $A_p$ is characteristic in $A$ and each $B_p$ is characteristic in $B$. Since $A,B\trianglelefteq G$, it follows that $A_p,B_p\trianglelefteq G$.

If $p\neq q$, then
\[
    [A_p,B_q]\leq A_p\cap B_q=1,
\]
because the commutator subgroup lies in both normal subgroups $A_p$ and $B_q$, while $A_p\cap B_q$ is both a $p$-group and a $q$-group. Thus different-prime factors commute. For each prime $p$, the product
\[
    A_pB_p
\]
is a normal $p$-subgroup of $AB$. Moreover,
\[
    AB=\prod_p A_pB_p,
\]
and the factors for distinct primes commute and intersect trivially. Therefore the Sylow subgroups of $AB$ are precisely the subgroups $A_pB_p$, and they are normal in $AB$. By Theorem~\ref{thm:8.7}, $AB$ is nilpotent.

Now let $\mathcal{N}$ be the set of nilpotent normal subgroups of $G$, and define
\[
    F=\prod_{N\in\mathcal{N}}N.
\]
Since $G$ is finite, it has only finitely many subgroups, so this is a finite product. Repeatedly applying the previous paragraph shows that $F$ is nilpotent. Also, a product of normal subgroups is normal, so $F\trianglelefteq G$. Finally, by construction, every nilpotent normal subgroup $N$ of $G$ is one of the factors in this product, and hence $N\leq F$. Thus $F$ is the largest nilpotent normal subgroup of $G$.
\end{solution}

\noindent The intersection of all maximal subgroups of a finite group $G$ is called the \emph{Frattini subgroup} of $G$ and is denoted by $\Phi(G)$.

\begin{exercise}
\label{ex:8.8}
Show that $\Phi(G)$ is a nilpotent normal subgroup of $G$.
\end{exercise}
\begin{solution}
If $G=1$, then the result is clear, so assume $G\neq 1$. First we show that $\Phi(G)\trianglelefteq G$. Every automorphism of $G$ sends maximal subgroups to maximal subgroups. Indeed, if $\alpha\in\Aut(G)$ and $M$ is maximal in $G$, then any subgroup strictly between $\alpha(M)$ and $G$ would pull back under $\alpha^{-1}$ to a subgroup strictly between $M$ and $G$. Thus automorphisms permute the maximal subgroups of $G$, and so preserve their intersection. Hence $\Phi(G)$ is characteristic in $G$, and therefore $\Phi(G)\trianglelefteq G$.

Now let $P$ be a Sylow $p$-subgroup of $\Phi(G)$. Since $\Phi(G)\trianglelefteq G$, the \nameref{thm:frattini-arg} gives
\[
    G=N_G(P)\Phi(G).
\]
If $N_G(P)<G$, then, since $G$ is finite, there is a maximal subgroup $M$ of $G$ with
\[
    N_G(P)\leq M<G.
\]
But $\Phi(G)$ is contained in every maximal subgroup of $G$, so $\Phi(G)\leq M$. Hence
\[
    G=N_G(P)\Phi(G)\leq M,
\]
a contradiction. Therefore $N_G(P)=G$, so $P\trianglelefteq G$, and in particular $P\trianglelefteq \Phi(G)$.

Thus every Sylow subgroup of $\Phi(G)$ is normal in $\Phi(G)$. By Theorem~\ref{thm:8.7}, $\Phi(G)$ is nilpotent.
\end{solution}

\begin{exercise}
\label{ex:8.9}
Show that $g \in \Phi(G)$ iff whenever $G = \langle S \rangle$ and $g \in S$, then $G = \langle S - \{g\} \rangle$.
\end{exercise}
\begin{solution}
First suppose that $g\in\Phi(G)$. Let $S$ be a subset of $G$ such that $G=\langle S\rangle$ and $g\in S$. Put
\[
    H=\langle S-\{g\}\rangle.
\]
We claim that $H=G$. If not, then $H$ is contained in some maximal subgroup $M$ of $G$. Since $g\in\Phi(G)$, and $\Phi(G)$ is contained in every maximal subgroup of $G$, we have $g\in M$. Also $S-\{g\}\subseteq H\leq M$. Hence every element of $S$ lies in $M$, so
\[
    G=\langle S\rangle\leq M<G,
\]
a contradiction. Therefore $G=\langle S-\{g\}\rangle$.

Conversely, suppose that $g\notin\Phi(G)$. Since $\Phi(G)$ is the intersection of all maximal subgroups of $G$, there is a maximal subgroup $M$ of $G$ such that $g\notin M$. Let
\[
    S=M\cup\{g\}.
\]
Then $\langle S\rangle=\langle M,g\rangle$. This subgroup contains $M$ properly, because $g\notin M$, and so maximality of $M$ forces
\[
    \langle S\rangle=G.
\]
However,
\[
    \langle S-\{g\}\rangle=\langle M\rangle=M\neq G.
\]
Thus $g$ does not have the stated removable-generator property. This proves the equivalence.
\end{solution}

\begin{exercise}
\label{ex:8.10}
Show that if $P$ is a finite $p$-group, then $P/\Phi(P)$ is an elementary abelian $p$-group.
\end{exercise}
\begin{solution}
If $P=1$, then the result is clear. Assume therefore that $P\neq 1$. Let $M$ be a maximal subgroup of $P$. By Proposition~\ref{prop:8.3} and Corollary~\ref{cor:8.4}, we have $M\trianglelefteq P$ and $\lvert P:M\rvert=p$. Hence $P/M$ is cyclic of order $p$.

Now let $x,y\in P$. Since $P/M$ has order $p$, the coset $xM$ has order dividing $p$, so
\[
    x^p\in M.
\]
Also $P/M$ is abelian, so
\[
    [x,y]\in M.
\]
These conclusions hold for every maximal subgroup $M$ of $P$. Therefore
\[
    x^p\in\Phi(P)
    \qquad\text{and}\qquad
    [x,y]\in\Phi(P)
\]
for all $x,y\in P$.

It follows that every element of $P/\Phi(P)$ has order dividing $p$, since
\[
    (x\Phi(P))^p=x^p\Phi(P)=\Phi(P),
\]
and all commutators are trivial in $P/\Phi(P)$, since
\[
    [x\Phi(P),y\Phi(P)]=[x,y]\Phi(P)=\Phi(P).
\]
Thus $P/\Phi(P)$ is an abelian $p$-group of exponent $p$. By Theorem~\ref{thm:8.6}, it is a direct product of cyclic $p$-groups; since every element has order dividing $p$, each non-trivial cyclic factor has order $p$. Therefore $P/\Phi(P)$ is elementary abelian.
\end{solution}

\end{document}
